% =====================================================================
% Interface e Explicabilidade
% Responsável: Pessoa 6
%
% Implementação em app.py. Falta inserir uma captura de tela da
% interface (ver \pendente abaixo).
% =====================================================================

\section{Interface e Explicabilidade}
\label{sec:interface}

A interface foi construída com Streamlit (\texttt{app.py}) e tem como objetivo tornar visível de onde vem cada parte da resposta. O usuário digita uma pergunta em linguagem natural e o sistema executa a busca híbrida (Seção~\ref{sec:recuperacao}), recupera os fatos e gera a resposta ancorada (Seção~\ref{sec:grounding}). A tela é organizada em cinco blocos:

\begin{enumerate}
  \item \textbf{Resposta}, com as citações \texttt{[F1]}, \texttt{[F2]}, \ldots{} no texto. Quando o sistema não encontra informação suficiente, a abstenção é destacada como aviso, e não como erro.
  \item \textbf{Indicadores}: o índice de confiança da recuperação, quantos dos fatos recebidos foram citados e se a fundamentação foi verificada. Se o LLM citar um fato inexistente, a interface informa que a citação foi removida.
  \item \textbf{Evidências}: para cada fato citado, a relação (origem, tipo e destino), o trecho da Wikipédia que a comprova, o documento de origem com link para a página, a confiança da extração e os identificadores dos trechos.
  \item \textbf{Subgrafo}: desenho, gerado com Graphviz, das entidades e relações enviadas ao LLM. As relações citadas na resposta aparecem destacadas em azul, com o número do fato correspondente.
  \item \textbf{Detalhes}: a lista completa de fatos enviados ao LLM, com a similaridade de cada um com a pergunta, as entidades encontradas pela busca híbrida e a consulta Cypher gerada.
\end{enumerate}

Na barra lateral, o usuário escolhe o número de saltos da travessia e se a resposta deve ser gerada pelo LLM. Sem um LLM configurado, a interface continua funcional no modo extrativo, listando os fatos mais relevantes com suas citações. Essa escolha de projeto garante que a explicabilidade não dependa de um serviço externo: o caminho da pergunta até a evidência é sempre exibido.

\pendente{inserir captura de tela da interface respondendo a uma pergunta multi-hop, por exemplo ``Que empresa foi fundada por pessoas que trabalharam na Fairchild Semiconductor?''}
